% from ega1-5-fr.tex
\oldpage[I]{140}
  \item[\rm(iv)] \emph{Si $f:X\to Y$ est un $S$-morphisme possédant la
    propriété \textbf{P}, tout $S'$-morphisme $f_{(S')}$ obtenu par une
    extension $S'\to S$ du préschéma de base, possède la propriété
    \textbf{P}.}
  \item[\rm(v)] \emph{Si le composé $g\circ f$ de deux morphismes
    $f:X\to Y$, $g:Y\to Z$ possède la propriété \textbf{P}, et si $g$ est
    séparé, $f$ possède la propriété \textbf{P}.}
  \item[\rm(vi)] \emph{Si un morphisme $f:X\to Y$ possède la propriété
    \textbf{P}, il en est de même de $f_{\mathrm{red}}$
    (\hyperref[I.5.1.5-fr]{5.1.5}).}
\end{enumerate}

Dans ces conditions, si on suppose (i) et (ii) vérifiées, (iii) et (iv) sont
\emph{équivalentes}, et (v) et (vi) sont des \emph{conséquences} de (i), (ii)
et (iii).

La première assertion a déjà été démontrée
(\hyperref[I.3.5.1-fr]{3.5.1}). Considérons la factorisation
(\hyperref[I.5.3.13-fr]{5.3.13}) de $f$ en
$X\xrightarrow{\Gamma_f}X\times_Z Y\xrightarrow{p_2}Y$~; la relation
$p_2=(g\circ f)\times_Z 1_Y$ montre que si $g\circ f$ possède la propriété
\textbf{P}, il en est de même de $p_2$ en vertu de (iii)~; si $g$ est séparé,
$\Gamma_f$ est une immersion fermée
(\hyperref[I.5.4.3-fr]{5.4.3}), donc possède aussi la propriété \textbf{P} par
(i)~; enfin, en vertu de (ii), $f$ possède la propriété \textbf{P}.

Enfin, considérons le diagramme commutatif
\[
  \xymatrix{
    X_{\mathrm{red}}\ar[r]^{f_{\mathrm{red}}}\ar[d] &
    Y_{\mathrm{red}}\ar[d]\\
    X\ar[r]_f &
    Y
  }
\]
où les flèches verticales sont des immersions fermées
(\hyperref[I.5.1.5-fr]{5.1.5}), donc possèdent la propriété \textbf{P} par
(i). L'hypothèse que $f$ possède la propriété \textbf{P} entraîne donc par
(ii) que
$X_{\mathrm{red}}\xrightarrow{f_{\mathrm{red}}}Y_{\mathrm{red}}\to Y$
possède la propriété \textbf{P}~; enfin, comme une immersion fermée est
séparée (\hyperref[I.5.5.1-fr]{5.5.1, (i)}), $f_{\mathrm{red}}$ possède la
propriété \textbf{P} en vertu de (v).

On notera que si on considère les propositions~:
\begin{enumerate}
  \item[\rm(i')] \emph{Toute immersion possède la propriété \textbf{P}.}
  \item[\rm(v')] \emph{Si $g\circ f$ possède la propriété \textbf{P}, il en
    est de même de $f$~;}
\end{enumerate}
alors le raisonnement fait ci-dessus montre que (v') est conséquence de (i'),
(ii) et (iii).
\end{remark}

\begin{env}[5.5.13]
\label{I.5.5.13-fr}
On notera que (v) et (vi) sont encore conséquences de (i), (iii) et
\begin{enumerate}
  \item[\rm(ii')] \emph{Si $j:X\to Y$ est une immersion fermée et
    $g:Y\to Z$ un morphisme possédant la propriété \textbf{P}, alors
    $g\circ j$ possède la propriété \textbf{P}.}
\end{enumerate}
De même, (v') est conséquence de (i'), (iii) et
\begin{enumerate}
  \item[\rm(ii'')] \emph{Si $j:X\to Y$ est une immersion et $g:Y\to Z$ un
    morphisme possédant la propriété \textbf{P}, alors $g\circ j$ possède la
    propriété \textbf{P}.}
\end{enumerate}
Cela résulte en effet aussitôt des raisonnements de
(\hyperref[I.5.5.12-fr]{5.5.12}).
\end{env}

% (continued in ega1-6-fr.tex, before its first page marker)
\section{Conditions de finitude}
\label{section:I.6-fr}

\subsection{Préschémas noethériens et localement noethériens.}
\label{subsection:I.6.1-fr}

\begin{definition}[6.1.1]
\label{I.6.1.1-fr}
On dit qu'un préschéma $X$ est \emph{noethérien} (resp. \emph{localement
noethérien}) s'il est réunion finie (resp. réunion) d'ouverts affines
$V_\alpha$ tels que l'anneau de chacun des schémas induits sur les $V_\alpha$
soit noethérien.
\end{definition}

Il résulte aussitôt de (\hyperref[I.1.5.2-fr]{1.5.2}) que si $X$ est
localement noethérien, le faisceau structural $\mathscr{O}_X$ est un
\emph{faisceau cohérent d'anneaux}, la question étant locale. \emph{Tout
sous-$\mathscr{O}_X$-}

